% !TEX root = ../main.tex
% Хавсралт А

\chapter{Судалгаанд ашигласан өгөгдлийн сангууд}
\label{AppendixA}
\pagecolor{white}

Энэхүү судалгаанд ясны хугарлын ангилал, сегментчлэл болон үнэлгээний даалгавруудад зориулан олон нийтэд нээлттэй хэд хэдэн өгөгдлийн санг ашиглав. MURA өгөгдлийг эхний хоёр туршилтын ангиллын суурь сургалтад ашигласан бол эцсийн модульчлагдсан pipeline-д бусад эх сурвалжуудыг ашигласан. Хүснэгт~\ref{tab:appA_datasets}-д ашигласан өгөгдлийн сангуудын үндсэн мэдээллийг харуулав.

\begin{table}[ht]
    \centering
    \begin{tabular}{|l|c|l|}
    \hline
    \textbf{Өгөгдлийн сан} & \textbf{Зургийн тоо} & \textbf{Даалгавар} \\ \hline
    MURA                  & $\sim$40{,}000 & Туршилт 1--2, ангиллын суурь сургалт \\ \hline
    FracAtlas             & 4{,}083  & Сегментчлэл, хугарлын зэрэг \\ \hline
    FracAtlas-COCO        & 966      & Туршилт 3-ын ангилал \\ \hline
    GRAZPEDWRI-DX         & 20{,}327 & Туршилт 3-ын ангилал \\ \hline
    YOLOBoneFracture      & 3{,}979  & Туршилт 3-ын ангилал, псевдо маск \\ \hline
    BoneFractureCVProject & 3{,}631  & Туршилт 3-ын ангилал, псевдо маск \\ \hline
    \end{tabular}
    \caption{Судалгаанд ашигласан өгөгдлийн сангууд}
    \label{tab:appA_datasets}
\end{table}

Эцсийн модульчлагдсан pipeline-д MURA-г оруулахгүйгээр нийтдээ 30{,}000 гаруй (нийлбэр $\approx$32{,}986) рентген зургийг боловсруулж ашигласан. Харин эхний хоёр туршилтын ангиллын суурь сургалтад MURA-ийн $\sim$40{,}000 зургийг тусад нь ашигласан болно.
